% ============================================================
% source.tex — Ch 01 : Arithmétique — Tronc Commun Maroc
% ============================================================

%% HEADER_APERCU — NE PAS EXTRAIRE
\documentclass[10pt, a4paper]{article}
% ==========================================
% 01_PACKAGES.TEX
% ==========================================
\usepackage[utf8]{inputenc}
\usepackage[T1]{fontenc}
\usepackage[french]{babel}
\usepackage[margin=1.5cm]{geometry} 
\usepackage{array} % DOIT ÊTRE CHARGÉ AVANT XCOLOR[TABLE] POUR ÉVITER LE BUG \insert@pcolumn
\usepackage[dvipsnames,table]{xcolor}
\usepackage{booktabs} % Professional tables
\usepackage{cellspace} % Espacement dynamique des cellules
\setlength{\cellspacetoplimit}{4pt}
\setlength{\cellspacebottomlimit}{4pt}
\usepackage{amsmath, amssymb, mathrsfs}
\usepackage[most]{tcolorbox}
\tcbuselibrary{xparse}
\usepackage{enumitem} 
\usepackage{fancyhdr}
\usepackage{tikz}
\usepackage{tkz-tab}

\usetikzlibrary{babel, shapes.geometric}
\usepackage{pgfplots}
\pgfplotsset{compat=newest}

\usepackage{multicol}
\usepackage{lmodern}
\usepackage{fourier}


\usepackage{palatino} % Main serif font
\usepackage[scaled]{helvet} % Professional Sans-Serif font (Helvetica)
\usepackage{microtype} % Improved typography
\usepackage{lipsum}

\usepackage{fontawesome5} % Pour les icônes modernes
\usepackage{longtable}
\usepackage{titlesec} % Contrôle absolu des espaces autour des titres
\usepackage{parskip}  % Supprime l'indentation et gère l'espacement des paragraphes

\usepackage{graphicx} %To insert the image properly without distorting its aspect ratio

\usepackage{setspace}

% Choisissez l'une des commandes suivantes :
%\onehalfspacing % Pour un interligne de 1.5
% \doublespacing % Pour un interligne double
 \setstretch{1.2} % Pour un espacement personnalisé très précis


% --- COLONNE ROBUSTE POUR CONTENU (Évite les bugs avec \section et tcolorbox) ---
\newcolumntype{M}[1]{>{\begin{minipage}[t]{\hsize}\vspace{1mm}}p{#1}<{\vspace{1mm}\end{minipage}}}


% ==========================================
% 02_STYLE.TEX
% ==========================================

% --- COULEURS ---
\definecolor{PrimaryColor}{HTML}{004DA0} % Thomas Dark Blue (Headings, Theorem Borders, primary elements)
\definecolor{SecondaryColor}{HTML}{0078D7} % Thomas Bright Blue (Main Graph lines, Highlights)
\definecolor{AccentColor}{HTML}{E31837} % Thomas Red (Asymptotes, important points)
\definecolor{NeutralColor}{HTML}{EBF3F9} % Soft Light Blue (Table headers)
\definecolor{LightBlue}{HTML}{EAF2F8}
\definecolor{BgColor}{HTML}{FAFAFA}
\definecolor{GraphGreen}{HTML}{1E8E3E} % <--- New Green added here



% --- GEOMETRY (Ajustement pour le cadre) ---
\geometry{a4paper, left=2.5cm, right=1.5cm, top=1.5cm, bottom=2cm}

% --- FORMAT ET ESPACEMENT DES TITRES (Hanging Numbers Style) ---
\titlespacing*{\section}{0pt}{4ex plus 1ex minus 0.5ex}{1.5ex plus 0.3ex}
\titlespacing*{\subsection}{0pt}{3ex plus 1ex minus 0.5ex}{1ex plus 0.2ex}
\titlespacing*{\subsubsection}{0pt}{2ex plus 0.5ex minus 0.2ex}{0.5ex plus 0.2ex}

\titleformat{\section}[hang]
  {\normalfont\sffamily\Large\bfseries\color{PrimaryColor}}
  {\llap{\colorbox{PrimaryColor}{\makebox[1cm][c]{\color{white}\sffamily\thesection}}\hspace{0.5cm}}}
  {0pt}
  {}

\titleformat{\subsection}[hang]
  {\normalfont\sffamily\large\bfseries\color{PrimaryColor}}
  {\llap{\makebox[1.5cm][l]{\color{PrimaryColor}\sffamily\thesubsection}}}
  {0pt}
  {}

\titleformat{\subsubsection}[hang]
  {\normalfont\sffamily\normalsize\bfseries\color{darkgray}}
  {\llap{\makebox[1.5cm][l]{\color{darkgray}\sffamily\thesubsubsection}}}
  {0pt}
  {}

% --- STYLE DES LISTES ET TABLEAUX ---

% --- GEOMETRIC ENUMERATE BADGES ---
\DeclareRobustCommand{\enumrect}[1]{\tcbox[on line, boxsep=0pt, boxrule=0pt, colback=PrimaryColor!80, coltext=white, sharp corners, fontupper=\bfseries\sffamily\footnotesize, left=3.5pt, right=3.5pt, top=2pt, bottom=2pt, box align=base, nobeforeafter]{#1}}

\DeclareRobustCommand{\enumcirc}[1]{%
  \tikz[baseline=(char.base)]{%
    \node[shape=circle, fill=PrimaryColor!70, text=white, font=\bfseries\sffamily\footnotesize, inner sep=1pt, minimum size=1em] (char) {#1};%
  }%
}

% Enumerate: 3 niveaux
\setlist[enumerate,1]{label=\textbf{\sffamily \arabic*.}, itemsep=4pt, topsep=4pt}
\setlist[enumerate,2]{label=\textbf{\textcolor{PrimaryColor}{\sffamily \alph*)}}, itemsep=2pt, topsep=2pt}
\setlist[enumerate,3]{label=\textbf{\textcolor{PrimaryColor!80}{\rmfamily\alph{enumii}\textsubscript{\arabic*})}}, itemsep=1pt, topsep=1pt}

% Itemize: 3 niveaux
\setlist[itemize,1]{label=\textcolor{black}{\textbullet}, itemsep=2pt, topsep=2pt}
\setlist[itemize,2]{label=\textcolor{black}{$\circ$}, itemsep=1pt, topsep=1pt}
\setlist[itemize,3]{label=\textcolor{black}{-}, itemsep=1pt, topsep=1pt}

% Tableaux professionnels
% rowcolors activé uniquement dans tabular (pas dans array/cases) via etoolbox
\usepackage{etoolbox}
\AtBeginEnvironment{tabular}{\rowcolors{2}{PrimaryColor!4}{white}}

% Utilisation de \AfterEndEnvironment pour s'exécuter HORS du tableau
% Utilisation de \hiderowcolors (commande officielle) pour stopper l'alternance
% \AfterEndEnvironment{tabular}{\hiderowcolors}  

\newcommand{\theadrow}{\rowcolor{NeutralColor}\color{PrimaryColor}\bfseries\sffamily}

% --- STYLE PGFPLOTS ---
\pgfplotsset{
    proplot/.style={
        axis lines=middle,
        axis line style={thick, black, -stealth},
        grid=none,
        tick label style={font=\sffamily\footnotesize, color=black},
        every axis plot/.append style={very thick, SecondaryColor},
        label style={font=\sffamily\small\itshape, black}
    }
}

% --- STYLE DES BOÎTES (Fiche Technique) ---
\tcbset{
    fichebox/.style={
        enhanced,
        boxrule=1pt,
        colframe=PrimaryColor,
        colback=white,
        coltitle=PrimaryColor,
        fonttitle=\normalfont\sffamily,
        attach boxed title to top left={xshift=5mm, yshift=-\tcboxedtitleheight/2},
        boxed title style={
            enhanced,
            colback=white, 
            frame hidden,   % Supprime la ligne parasite
            boxrule=0pt,
            arc=3pt,
            left=2mm, right=2mm, top=1mm, bottom=1mm
        },
        top=4mm, bottom=2mm, left=3mm, right=3mm
    }
}

% --- STYLE "COURS" (Pages avec cadre) ---
\fancypagestyle{cours}{
    \fancyhf{} 
    \renewcommand{\headrulewidth}{0pt} 
    \renewcommand{\footrulewidth}{0pt} 
    
    \fancyfoot[L]{\small Ch 11. Fonctions numériques} 
    \fancyfoot[C]{\small \thepage}                    
    \fancyfoot[R]{\small Pr. Kaazouzi Ayoub}         
}

% ==========================================
% 02_STYLE.TEX : STYLE ATTACHED BOXED TITLE SANS ICÔNES (Sans Numérotation)
% ==========================================

% --- PALETTE DE COULEURS DOUCES (SOFT COLORS) ---
\definecolor{CAct}{HTML}{67899F}    % Activité (Soft Steel Blue)
\definecolor{CDef}{HTML}{C16B67}    % Définition (Soft Rose/Coral)
\definecolor{CVoc}{HTML}{8477A4}    % Vocabulaire (Soft Lavender)
\definecolor{CProp}{HTML}{C59145}   % Propriété (Soft Gold/Amber)
\definecolor{CPropo}{HTML}{6C8EA8}  % Proposition (Soft Blue)
\definecolor{CThm}{HTML}{729C7A}    % Théorème (Soft Sage Green)
\definecolor{CCor}{HTML}{619489}    % Corollaire (Soft Mint/Teal)
\definecolor{CPre}{HTML}{8F857B}    % Preuve (Soft Taupe)
\definecolor{CEx}{HTML}{C87D53}     % Exemple (Soft Earthy Orange)
\definecolor{CRem}{HTML}{8895A5}    % Remarque (Soft Cool Gray)
\definecolor{CApp}{HTML}{0078D7}    % Application (Thomas Bright Blue)
\definecolor{CExer}{HTML}{004DA0}   % Exercice (Thomas Dark Blue)
\definecolor{CTech}{HTML}{8B1A1A}   % Techniques (Dark Red / Burgundy)
\definecolor{CSol}{HTML}{4A8C5C}    % Solution (Soft Forest Green)

% --- STYLE DE BASE DES BOÎTES (Option B Modifiée - 3 param) ---
\tcbset{
    myboxstyle/.style n args={3}{
        enhanced,
        colback=white,
        colframe=#1,
        boxrule=1pt,
        arc=4pt,
        coltitle=white,
        fonttitle=\normalfont\sffamily\bfseries,
        attach boxed title to top left={xshift=4mm, yshift=-\tcboxedtitleheight/2},
        boxed title style={
            enhanced,
            colback=#1,
            colframe=#1,
            arc=3pt,
            boxrule=0pt,
        },
        title={#2},
        overlay unbroken and first={
            \ifstrempty{#3}{}{
                \node[
                    fill=#1,
                    text=white,
                    font=\bfseries\small,
                    rounded corners=2.5pt,
                    inner xsep=2mm,
                    inner ysep=1.5mm,
                    anchor=east
                ] at ([xshift=-4mm]frame.north east) {\faClock~#3};
            }
        },
        top=3mm, bottom=3mm, left=4mm, right=4mm
    },
    emptyboxstyle/.style n args={3}{
        enhanced,
        frame hidden,
        interior hidden,
        boxrule=0pt,
        coltitle=#1,
        fonttitle=\normalfont\sffamily\bfseries,
        title={\rule[0.05em]{0.7cm}{0.6em}\hspace{2mm}#2},
        titlerule=0pt,
        boxsep=0pt,
        left=0pt, right=0pt, top=2mm, bottom=2mm,
        toptitle=2mm, bottomtitle=2mm,
        overlay unbroken and first={
            \ifstrempty{#3}{}{
                \node[
                    text=#1,
                    font=\bfseries\small,
                    anchor=east
                ] at ([yshift=-3.5mm]frame.north east) {\faClock~#3};
            }
        }
    },
    leftbarstyle/.style n args={3}{
        enhanced,
        frame hidden,
        interior hidden,
        boxrule=0pt,
        borderline west={3pt}{0pt}{#1},
        coltitle=#1,
        fonttitle=\normalfont\sffamily\bfseries,
        title={#2},
        titlerule=0pt,
        boxsep=0pt,
        left=4mm, right=0mm, top=2mm, bottom=2mm,
        toptitle=2mm, bottomtitle=2mm,
        overlay unbroken and first={
            \ifstrempty{#3}{}{
                \node[
                    text=#1,
                    font=\bfseries\small,
                    anchor=east
                ] at ([yshift=-3.5mm]frame.north east) {\faClock~#3};
            }
        }
    },
    topbottomstyle/.style n args={3}{
        enhanced,
        frame hidden,
        colback=white,
        colbacktitle=#1!15,
        coltitle=#1,
        fonttitle=\normalfont\sffamily\large,
        title={#2},
        boxrule=0pt,
        titlerule=0pt,
        borderline south={1pt}{0pt}{#1},
        left=4mm, right=4mm, top=3mm, bottom=3mm,
        toptitle=2.5mm, bottomtitle=2.5mm,
        overlay unbroken and first={
            \ifstrempty{#3}{}{
                \node[
                    text=#1,
                    font=\bfseries\small,
                    anchor=east
                ] at ([xshift=-4mm, yshift=-4mm]frame.north east) {\faClock~#3};
            }
        }
    },
    topstyle/.style n args={3}{
        enhanced,
        frame hidden,
        colback=white,
        colbacktitle=#1!15,
        coltitle=#1,
        fonttitle=\normalfont\sffamily\large,
        title={#2},
        boxrule=0pt,
        titlerule=0pt,
        left=4mm, right=4mm, top=3mm, bottom=3mm,
        toptitle=2.5mm, bottomtitle=2.5mm,
        overlay unbroken and first={
            \ifstrempty{#3}{}{
                \node[
                    text=#1,
                    font=\bfseries\small,
                    anchor=east
                ] at ([xshift=-4mm, yshift=-4mm]frame.north east) {\faClock~#3};
            }
        }
    },
    remarkstyle/.style n args={3}{
        enhanced,
        frame hidden,
        colback=#1!10,
        boxrule=0pt,
        coltitle=#1,
        fonttitle=\normalfont\sffamily\bfseries,
        attach title to upper,
        after title={\enskip},
        title={#2},
        left=3mm, right=3mm, top=2mm, bottom=2mm,
        overlay unbroken and first={
            \ifstrempty{#3}{}{
                \node[
                    text=#1,
                    font=\bfseries\small,
                    anchor=east
                ] at ([yshift=-3.5mm]frame.north east) {\faClock~#3};
            }
        }
    },
    justtitlestyle/.style n args={3}{
        enhanced,
        frame hidden,
        colback=white,
        coltitle=#1,
        fonttitle=\normalfont\sffamily\Large\bfseries,
        title={#2},
        boxrule=0pt,
        titlerule=0pt,
        boxsep=0pt,
        left=0mm, right=0mm, top=2mm, bottom=2mm,
        toptitle=2mm, bottomtitle=2mm,
        overlay unbroken and first={
            \ifstrempty{#3}{}{
                \node[
                    text=#1,
                    font=\bfseries\small,
                    anchor=east
                ] at ([yshift=-3.8mm]frame.north east) {\faClock~#3};
            }
        }
    },
    techniquestyle/.style n args={3}{
        enhanced,
        frame hidden,
        colback=white,
        coltitle=#1,
        fonttitle=\normalfont\rmfamily\scshape\Large\bfseries,
        title={#2},
        halign title=center,
        boxrule=0pt,
        titlerule=0pt,
        boxsep=0pt,
        left=0mm, right=0mm, top=2mm, bottom=2mm,
        toptitle=4mm, bottomtitle=2mm,
        overlay unbroken and first={
            \ifstrempty{#3}{}{
                \node[
                    text=#1,
                    font=\bfseries\small,
                    anchor=east
                ] at ([yshift=-5.8mm]frame.north east) {\faClock~#3};
            }
        }
    },
    inlinestyle/.style n args={3}{
        enhanced,
        colback=white,
        colframe=#1,
        boxrule=1pt,
        sharp corners,
        coltitle=#1,
        fonttitle=\normalfont\sffamily\bfseries,
        attach title to upper,
        after title={\enskip},
        title={#2},
        left=3mm, right=3mm, top=2mm, bottom=2mm,
        overlay unbroken and first={
            \ifstrempty{#3}{}{
                \node[
                    text=#1,
                    font=\bfseries\small,
                    anchor=east
                ] at ([yshift=-3.5mm]frame.north east) {\faClock~#3};
            }
        }
    }
}

% ==========================================
% ENVIRONNEMENTS ENCADRÉS (SANS NUMÉROTATION)
% ==========================================

% 1. Activité
\NewTColorBox[auto counter]{Activite}{ O{} O{} }{
    emptyboxstyle={CAct}{ACTIVITÉ \thetcbcounter\ifstrempty{#1}{}{~(#1)}}{#2},
    breakable
}

% 2. Définition
\NewTColorBox{Definition}{ O{} O{} }{
    inlinestyle={CDef}{DÉFINITION\ifstrempty{#1}{}{~(#1)}}{#2}
}

% 3. Vocabulaire  (Gris / Noir)
\NewTColorBox[auto counter]{Vocabulaire}{ O{} O{} }{
    emptyboxstyle={black!70}{VOCABULAIRE \ifstrempty{#1}{}{~(#1)}}{#2}
}

% 3. Notations (Gris / Noir)
\NewTColorBox[auto counter]{Notations}{ O{} O{} }{
    emptyboxstyle={black!70}{NOTATIONS\ifstrempty{#1}{}{~(#1)}}{#2}
}

% 4. Propriété
\NewTColorBox{Propriete}{ O{} O{} }{
    inlinestyle={CProp}{PROPRIÉTÉ\ifstrempty{#1}{}{~(#1)}}{#2}
}

% 5. Proposition
\NewTColorBox{Proposition}{ O{} O{} }{
    inlinestyle={CPropo}{PROPOSITION\ifstrempty{#1}{}{~(#1)}}{#2}
}

% 6. Théorème
\NewTColorBox{Theoreme}{ O{} O{} }{
    inlinestyle={CThm}{THÉORÈME\ifstrempty{#1}{}{~(#1)}}{#2}
}

% 7. Corollaire
\NewTColorBox{Corollaire}{ O{} O{} }{
    inlinestyle={CCor}{COROLLAIRE\ifstrempty{#1}{}{~(#1)}}{#2}
}

% 8. Preuve
\NewTColorBox{Preuve}{ O{} O{} }{
    emptyboxstyle={CPre}{PREUVE\ifstrempty{#1}{}{~(#1)}}{#2},
    after upper={\hfill$\blacksquare$}
}

% 9. Exemple
\NewTColorBox{Exemple}{ O{} O{} }{
    emptyboxstyle={CEx}{EXEMPLES\ifstrempty{#1}{}{~(#1)}}{#2}
}

% 10. Remarques
\NewTColorBox{Remarque}{ O{} O{} }{
    remarkstyle={CRem}{\faInfoCircle\hspace{1mm}Remarque\ifstrempty{#1}{}{~(#1)}}{#2}
}

% 11. Application
\NewTColorBox[auto counter]{Application}{ O{} O{} }{
    emptyboxstyle={CApp}{APPLICATION \thetcbcounter\ifstrempty{#1}{}{~(#1)}}{#2},
    breakable
}

% 12. Exercice
\NewTColorBox[auto counter]{Exercice}{ O{} O{} }{
    emptyboxstyle={CExer}{EXERCICE DE SYNTHÈSE \thetcbcounter\ifstrempty{#1}{}{~(#1)}}{#2},
    breakable
}

% 13. Techniques
\NewTColorBox{Techniques}{ O{} O{} }{
    techniquestyle={CTech}{Techniques\ifstrempty{#1}{}{~(#1)}}{#2}
}

% ==========================================
% ENVIRONNEMENTS TEXTUELS SIMPLES RESTANTS
% ==========================================

% 14. Solution
\NewTColorBox{Solution}{ O{} O{} }{
    emptyboxstyle={CSol}{SOLUTION\ifstrempty{#1}{}{~(#1)}}{#2},
    breakable
}
% ==========================================
% 03_MACROS.TEX : Custom Math Shortcuts
% ==========================================

\newcommand{\N}{\mathbb{N}}
\newcommand{\Z}{\mathbb{Z}}
\newcommand{\R}{\mathbb{R}}
\newcommand{\C}{\mathbb{C}}

% --- Geometry & Vectors ---
\newcommand{\vect}[1]{\overrightarrow{#1}}
\newcommand{\norm}[1]{\left\lVert#1\right\rVert}
\newcommand{\abs}[1]{\left\lvert#1\right\rvert}

% --- Analysis ---
\newcommand{\dd}{\mathop{}\!\mathrm{d}} % Correct upright 'd' for integrals (\dd x)
\newcommand{\e}{\mathrm{e}}             % Upright 'e' for the exponential function
\newcommand{\limite}[2]{\lim_{#1 \to #2}}

% --- Toujours utiliser \dfrac (grande fraction) ---
\let\oldfrac\frac
\renewcommand{\frac}[2]{\dfrac{#1}{#2}}

% --- Utility ---
\newcommand{\ds}{\displaystyle} % Forces large fractions and limits inline
\newcommand{\duree}[1]{\hfill\small\faClock~#1}
\geometry{a4paper, left=2.2cm, right=0.6cm, top=0.8cm, bottom=1.5cm}

\begin{document}
\pagestyle{cours}
%% FIN_HEADER_APERCU

%% SECTION
\section{Ensemble $\N$ --- Nombres pairs et nombres impairs}

%% DEFINITION
\begin{Definition}[Notations et vocabulaires]
Les nombres $0 ; 1 ; 2 ; 3 ; \ldots$ forment un ensemble appelé \textbf{ensemble des entiers naturels}, noté $\N$ :
$$\N = \{0\,;\,1\,;\,2\,;\,3\,;\,4\,;\,5\,;\,\ldots\,;\,99\,;\,\ldots\}$$
On note $\N^* = \N \setminus \{0\}$ l'ensemble des entiers naturels \textbf{non nuls} ($\N^*$ se lit « $\N$ privé de zéro »).
\end{Definition}

%% DEFINITION
\begin{Definition}[Nombres pairs et nombres impairs]
Soit $a$ un nombre entier naturel.
\begin{itemize}
  \item Un \textbf{nombre pair} est un multiple de $2$.
  \item Un \textbf{nombre impair} est un nombre qui n'est pas pair.
\end{itemize}
\end{Definition}

%% REMARQUE
\begin{Remarque}
\begin{itemize}
  \item Il existe une infinité d'entiers naturels.
  \item Un nombre pair se termine nécessairement par $0$, $2$, $4$, $6$ ou $8$.
  \item Un nombre impair se termine nécessairement par $1$, $3$, $5$, $7$ ou $9$.
\end{itemize}
\end{Remarque}

%% PROPRIETE
\begin{Propriete}[Écriture d'un nombre pair ou impair quelconque]
Un nombre naturel $n$ est \textbf{pair} si et seulement s'il s'écrit :
$$n = 2k \quad \text{avec } k \in \N$$
Un nombre naturel $n$ est \textbf{impair} si et seulement s'il s'écrit :
$$n = 2k + 1 \quad \text{avec } k \in \N$$
\end{Propriete}

%% PROPRIETE
\begin{Propriete}[Opérations sur les nombres pairs et impairs]
\begin{itemize}
  \item La somme de deux nombres de \textbf{même parité} est \textbf{paire}.
  \item La somme de deux nombres de \textbf{parité différente} est \textbf{impaire}.
  \item Si l'un des facteurs d'un produit est pair, alors ce produit est \textbf{pair}.
  \item Un nombre élevé à une puissance conserve sa parité.
\end{itemize}
\end{Propriete}

%% PROPRIETE
\begin{Propriete}[Nombres consécutifs]
Soient $n$ et $n+1$ deux entiers consécutifs.
\begin{itemize}
  \item La somme de deux entiers consécutifs est \textbf{impaire}.
  \item Le produit de deux entiers consécutifs est \textbf{pair}.
\end{itemize}
\end{Propriete}

%% SECTION
\section{Multiples et diviseurs d'un entier naturel}

%% DEFINITION
\begin{Definition}[Multiples et diviseurs]
Soient $a$ et $b$ deux entiers naturels. S'il existe un entier naturel $k$ tel que $a = k \times b$, on dit que :
\begin{itemize}
  \item $b$ \textbf{divise} $a$ (ou $b$ est un \textbf{diviseur} de $a$),
  \item $a$ est \textbf{divisible} par $b$,
  \item $a$ est un \textbf{multiple} de $b$.
\end{itemize}
\end{Definition}

%% REMARQUE
\begin{Remarque}
\begin{itemize}
  \item $0$ est un multiple de tous les nombres entiers naturels.
  \item Tout entier naturel est divisible par $1$ et par lui-même.
\end{itemize}
\end{Remarque}

%% PROPRIETE
\begin{Propriete}[Opérations sur les multiples]
Soient $a$, $b$ et $c$ des entiers naturels.
\begin{itemize}
  \item Si $a$ est multiple de $c$ et $b$ est multiple de $c$, alors $a + b$ est multiple de $c$ ;
        et si de plus $a \geq b$, alors $a - b$ est multiple de $c$.
  \item Si $a$ est multiple de $b$ et $b$ est multiple de $c$, alors $a$ est multiple de $c$ \emph{(transitivité)}.
  \item Si $c$ est multiple de $a$, alors $nc$ est multiple de $a$ pour tout $n \in \N$.
\end{itemize}
\end{Propriete}

%% PROPRIETE
\begin{Propriete}[Opérations sur les diviseurs]
Soient $a$, $b$ et $c$ des entiers naturels.
\begin{itemize}
  \item Si $a$ divise $b$ et $c$, et $b \geq c$, alors $a$ divise $b + c$ et $b - c$.
  \item Si $a$ divise $b$, alors $a$ divise $bc$.
\end{itemize}
\end{Propriete}

%% PROPRIETE
\begin{Propriete}[Critères de divisibilité]
\begin{itemize}
  \item Divisible par $2$ : le chiffre des unités est pair.
  \item Divisible par $3$ : la somme des chiffres est divisible par $3$.
  \item Divisible par $4$ : le nombre formé par les deux derniers chiffres est divisible par $4$.
  \item Divisible par $5$ : se termine par $0$ ou $5$.
  \item Divisible par $6$ : divisible à la fois par $2$ et par $3$.
  \item Divisible par $9$ : la somme des chiffres est divisible par $9$.
\end{itemize}
\end{Propriete}

%% SECTION
\section{Les nombres premiers}

%% DEFINITION
\begin{Definition}[Nombre premier]
Un nombre est \textbf{premier} s'il possède exactement deux diviseurs : $1$ et lui-même.
\end{Definition}

%% REMARQUE
\begin{Remarque}
\begin{itemize}
  \item $1$ n'est \textbf{pas} un nombre premier : il n'a qu'un seul diviseur.
  \item $2$ est le seul nombre premier pair. Tous les autres nombres premiers sont impairs.
  \item Tout entier naturel $\geq 2$ admet au moins un diviseur premier.
  \item L'ensemble des nombres premiers est infini (théorème d'Euclide).
\end{itemize}
\end{Remarque}

%% PROPRIETE
\begin{Theoreme}[Test de primalité --- critère d'arrêt]
Pour qu'un entier $n \geq 2$ soit premier, il faut et il suffit qu'aucun nombre premier $p$
vérifiant $p \leq \sqrt{n}$ ne divise $n$.
\end{Theoreme}

%% PROPRIETE
\begin{Theoreme}[Décomposition en produit de facteurs premiers]
Tout entier naturel supérieur ou égal à $2$ se décompose de façon \textbf{unique} en produit
de facteurs premiers.
\end{Theoreme}

%% SECTION
\section{PGCD et PPCM}

%% DEFINITION
\begin{Definition}[PGCD et PPCM]
Soient $a$ et $b$ deux entiers naturels non nuls.
\begin{itemize}
  \item Un \textbf{diviseur commun} de $a$ et $b$ est un entier qui divise à la fois $a$ et $b$.
  \item Le \textbf{PGCD} (plus grand commun diviseur) se note $\operatorname{pgcd}(a,b)$ ou $a \wedge b$.
  \item Le \textbf{PPCM} (plus petit multiple commun) se note $\operatorname{ppcm}(a,b)$ ou $a \vee b$.
\end{itemize}
\end{Definition}

%% PROPRIETE
\begin{Propriete}[Calcul du PGCD et du PPCM par décomposition en facteurs premiers]
Soient $a$ et $b$ deux entiers naturels non nuls.
\begin{itemize}
  \item $\operatorname{pgcd}(a,b)$ est le produit des facteurs premiers \textbf{communs} élevés au \textbf{plus petit} exposant.
  \item $\operatorname{ppcm}(a,b)$ est le produit de \textbf{tous} les facteurs premiers (communs ou non) élevés au \textbf{plus grand} exposant.
\end{itemize}
\end{Propriete}

%% REMARQUE
\begin{Remarque}
Soient $a$ et $b$ deux entiers naturels non nuls.
\begin{itemize}
  \item $\operatorname{pgcd}(a,a) = a$ \quad;\quad $\operatorname{pgcd}(1,a) = 1$ \quad;\quad $\operatorname{pgcd}(a,0) = a$.
  \item $\operatorname{ppcm}(a,a) = a$ \quad;\quad $\operatorname{ppcm}(1,a) = a$ \quad;\quad $\operatorname{ppcm}(a,0) = 0$.
  \item Si $a$ divise $b$, alors $\operatorname{pgcd}(a,b) = a$ et $\operatorname{ppcm}(a,b) = b$.
\end{itemize}
\end{Remarque}

%% PROPRIETE
\begin{Propriete}[Relation entre PGCD et PPCM]
Soient $a$ et $b$ deux entiers naturels non nuls. On a :
$$\operatorname{pgcd}(a,b) \times \operatorname{ppcm}(a,b) = a \times b$$
\end{Propriete}

%% DEFINITION
\begin{Definition}[Nombres premiers entre eux]
Deux entiers $a$ et $b$ sont \textbf{premiers entre eux} si $\operatorname{pgcd}(a,b) = 1$,
c'est-à-dire s'ils n'ont aucun diviseur commun autre que $1$.
\end{Definition}

%% REMARQUE
\begin{Remarque}
Il ne faut pas confondre \textbf{nombre premier} et \textbf{nombres premiers entre eux}.
Par exemple, $15$ et $22$ sont premiers entre eux, mais ne sont pas des nombres premiers.
\end{Remarque}

%% TECHNIQUES
\begin{Techniques}[Algorithme d'Euclide --- Calcul du PGCD]
Pour calculer $\operatorname{pgcd}(a,b)$ avec $a > b$, on applique la règle :
$$\operatorname{pgcd}(a,b) = \operatorname{pgcd}(b,r)$$
où $r$ est le reste de la division euclidienne de $a$ par $b$ (i.e.\ $a = bq + r$, $r < b$).

On répète l'opération jusqu'à obtenir un reste nul. Le PGCD est le \textbf{dernier reste non nul}.

\textbf{Exemple :} Calculer $\operatorname{pgcd}(255,\,141)$.
\begin{align*}
  255 &= 1 \times 141 + 114 \\
  141 &= 1 \times 114 + 27  \\
  114 &= 4 \times 27  + 6   \\
  27  &= 4 \times 6   + 3   \\
  6   &= 2 \times 3   + 0
\end{align*}
Donc $\operatorname{pgcd}(255,\,141) = 3$.
\end{Techniques}

\end{document}
